\documentclass[10pt,a4paper]{article}

\usepackage[utf8]{inputenc}
\usepackage[T1]{fontenc}
\usepackage[spanish,es-noshorthands]{babel}
\usepackage[top=1.3cm,bottom=1.3cm,left=1.5cm,right=1.5cm]{geometry}
\usepackage{booktabs}
\usepackage{graphicx}
\usepackage{xcolor}
\usepackage[colorlinks=true,linkcolor=black,urlcolor=blue]{hyperref}
\usepackage{titlesec}
\usepackage{enumitem}
\usepackage{paralist}

\titleformat{\section}{\normalsize\bfseries}{\thesection.}{0.5em}{}
\titlespacing{\section}{0pt}{4pt}{2pt}
\setlist[itemize]{leftmargin=1.1em,itemsep=0pt,parsep=0pt,topsep=1pt}
\setlength{\parskip}{2pt}
\setlength{\parindent}{0pt}
\pagestyle{empty}

\begin{document}

{\large\bfseries Gestión de cobranza: humanos vs.\ agentes de IA}\hfill
{\small Andrés Chaparro · Sept.\ 2026}\\[-2pt]
\rule{\textwidth}{0.4pt}

%------------------------------------------------------------------
\section{Datos}

100 grabaciones de cobranza telefónica (50 humanas, 50 IA), WAV mono 8\,kHz/16\,bit,
\textbf{381 min} totales. Contenido homogéneo: negociación de una obligación financiera
(saldo, descuento, cuotas). Todas llevan un \textbf{pitido de censura} que enmascara datos
personales: protege privacidad, pero añade ruido y trunca texto.

\begin{table}[h]
\centering\small
\vspace{-4pt}
\begin{tabular}{lcccc}
\toprule
& \multicolumn{2}{c}{\textbf{Humanos} (n=50)} & \multicolumn{2}{c}{\textbf{IA} (n=50)} \\
\cmidrule(lr){2-3}\cmidrule(lr){4-5}
& Media & Mediana & Media & Mediana \\
\midrule
Duración (s)          & 238,1 & 171,1 & 219,3 & 149,0 \\
Segmentos             & 74,0  & 53,0  & 54,1  & 34,5 \\
Palabras              & 500,4 & 332,5 & 441,0 & 296,5 \\
Turnos/min\,$^{*}$    & 3,72  & ---   & 4,67  & --- \\
\bottomrule
\end{tabular}
\vspace{-2pt}
\caption*{\footnotesize $^{*}$Solo llamadas con dos hablantes bien detectados.}
\end{table}
\vspace{-10pt}

\begin{itemize}
  \item \textbf{Dispersión enorme:} de 60\,s a 1.233\,s; la desviación estándar es del
  orden de la media $\Rightarrow$ se usan medianas, no promedios.
  \item \textbf{Los humanos hablan más tiempo, la IA alterna más.} Con conversación
  efectiva, mediana humana \textbf{186\,s} vs.\ \textbf{136,5\,s} en IA; pero la IA
  registra \textbf{4,67 turnos/min} vs.\ \textbf{3,72}.
\end{itemize}

%------------------------------------------------------------------
\section{Calidad de la medición}

Los datos \textbf{no se observan, se estiman}: cada variable pasa por dos modelos
(transcripción y diarización) y cada uno aporta error. El diagnóstico está en los hablantes
detectados por llamada, que deberían ser siempre dos:

\begin{table}[h]
\centering\small
\vspace{-4pt}
\begin{tabular}{lccc}
\toprule
& 1 hablante & 2 hablantes & 3 etiquetas \\
\midrule
Humanos & \textbf{16} (32\,\%) & 33 (66\,\%) & 1 (2\,\%) \\
IA      & 1 (2\,\%)            & \textbf{46} (92\,\%) & 3 (6\,\%) \\
\bottomrule
\end{tabular}
\end{table}
\vspace{-12pt}

\begin{itemize}
  \item \textbf{Los 16 casos humanos de 1 hablante no son falta de contacto:} 15 son
  conversaciones reales cuya diarización colapsó ambas voces (ej.\ \texttt{1c72dd55}:
  57 segmentos de negociación, todos en una etiqueta). Solo \texttt{0c6ada7d} es buzón real.
  \item \textbf{Las ``3 etiquetas'' no son un tercer interlocutor:} son 2 hablantes más una
  etiqueta \texttt{UNKNOWN} residual, con $<$6\,s de audio en los cuatro casos.
  \item \textbf{El error se concentra en turnos $<$1,5\,s} (``sí, con ella'',
  ``bien, gracias''). Verificamos que falla la diarización misma, no su combinación con el
  texto: un fragmento de 0,6\,s genera un \textit{embedding} demasiado corto para separarse
  del hablante vecino. Limitación conocida del
  enfoque.\footnote{\href{https://docs.nvidia.com/nemo-framework/user-guide/24.07/nemotoolkit/asr/speaker_diarization/models.html}{NVIDIA NeMo}, diarización. Parafraseado.}
\end{itemize}

\textbf{Por qué la IA se diariza mejor (92\,\% vs.\ 66\,\%):} voz sintética de timbre y
volumen constantes, sin carga emocional $\Rightarrow$ \textit{embedding} compacto y muy
distinto al humano; en cambio dos humanos comparten canal de 8\,kHz, acento y rango vocal
similares; además la voz de la IA se inyecta digitalmente, sin micrófono ni ruido de
\textit{call center}.

\textbf{Consecuencia analítica:} toda métrica que dependa de \emph{quién} habla tiene sesgo
sistemático contra las llamadas humanas. Esas variables se reportan solo sobre el
subconjunto con dos hablantes; las basadas en texto agregado usan la muestra completa.

\textbf{Decisiones:} se excluyen del análisis de conducta el buzón humano
(\texttt{0c6ada7d}) y el no-contacto de IA (\texttt{edbc4154}) --- sin interacción que
evaluar, aunque sí cuentan como contactabilidad. \textbf{No} se excluyen las 15
conversaciones mal diarizadas: su texto es válido y descartarlas sesgaría la muestra hacia
las llamadas más limpias. Pitido de censura y error en turnos breves se asumen como
limitaciones declaradas.

%------------------------------------------------------------------
\section{Construcción de la base}

Cada \texttt{.wav} se convierte en un JSON con la conversación segmentada
(\texttt{inicio, fin, hablante, texto}). La unidad es el \textbf{turno}, no la llamada
agregada: los minutos totales por hablante no permiten medir alternancia, latencia ni
contexto (``qué respondió el agente \emph{después} de una objeción''). Todo corre en
\textbf{Google Colab} con GPU T4.

\textbf{Transcripción --- Whisper \texttt{large-v3}} vía \texttt{faster-whisper}
(\texttt{CTranslate2}, \texttt{float16}). Se descartó \texttt{whisperx} por dependencias
incompatibles con Colab. La elección de \texttt{large-v3} sobre \texttt{small} fue empírica:
en el mismo audio, \texttt{small} transcribió ``con dos besos y ocho pesos'' donde
\texttt{large-v3} recuperó ``a ver, dos, dos, tres, cuatro''.
\begin{itemize}
  \item \texttt{vad\_filter} con \texttt{threshold=0.2}: con el valor por defecto (0,5) se
  perdían los primeros 50\,s de llamadas con volumen bajo; desactivarlo producía
  alucinaciones (texto inventado en inglés sobre silencios).
  \item \texttt{condition\_on\_previous\_text=False}: evita que un error se propague.
  \texttt{no\_speech\_threshold=0.6} y \texttt{compression\_ratio\_threshold=2.4}: descartan
  tramos sin voz y texto repetitivo, síntomas típicos de alucinación.
\end{itemize}

\textbf{Diarización --- \texttt{pyannote.audio} 3.1.} Tres pasos: segmenta tramos de voz,
calcula un \textit{embedding} (huella vectorial) por tramo, y los \textbf{agrupa por
similitud} para resolver quién habla.
\begin{itemize}
  \item \texttt{min\_speakers=max\_speakers=2}: el dominio lo garantiza (agente + deudor).
  \item \texttt{min\_cluster\_size}: 12 (defecto) $\rightarrow$ \textbf{5}. Define cuántos
  \textit{embeddings} necesita un grupo para ser hablante propio; con 12, las intervenciones
  breves quedaban absorbidas por el hablante dominante. Con menos de 5 aparecen micro-turnos
  espurios sin ganar precisión.
  \item Cada segmento de Whisper recibe el hablante con \textbf{mayor solapamiento
  temporal}. Se probó asignación palabra por palabra (\texttt{word\_timestamps}): los
  \textit{timestamps} de palabra resultaron inestables --- turnos de duración cero y
  etiquetas vacías --- y se descartó.
\end{itemize}

\end{document}
